\subsection{INVERSE LIMITS OF TOPOLOGICAL GROUPS AND SPACES WITH OPERATORS}

We shall say that an inverse system \((\mathbf{G}_{\alpha}, f_{\alpha\beta})\) is an inverse system of topological groups if the \(G_{\alpha}\) are topological groups and the \(f_{\alpha\beta}\) are continuous homomorphisms. Then \(G = \varprojlim G_{\alpha}\) is a subgroup of the product group \(\prod_{\alpha} G_{\alpha}\) ; if we endow G with the topological group structure induced by that of \(\prod_{\alpha} G_{\alpha}\) , then the topological group so obtained is called the inverse limit of the inverse system of topological groups \((\mathbf{G}_{\alpha}, f_{\alpha\beta})\) . If the \(G_{\alpha}\) are Hausdorff (resp. Hausdorff and complete) then G is Hausdorff and closed in \(\prod_{\alpha} G_{\alpha}\) (resp. Hausdorff and complete) (Chapter I, § 8, no. 2, Proposition 7, Corollary 2 and Chapter II, § 3, no. 5, Corollary to Proposition 10).

If \((\mathbf{G}_{\alpha}^{\prime}, f_{\alpha\beta}^{\prime})\) is another inverse system of topological groups and if, for each \(\alpha\) , \(u_{\alpha}: G_{\alpha} \to G_{\alpha}^{\prime}\) is a continuous homomorphism such that the \(u_{\alpha}\) form an inverse system of mappings, then \(u = \varprojlim u_{\alpha}\) is a continuous homomorphism of G into \(G^{\prime} = \varprojlim G_{\alpha}^{\prime}\) (Chapter I, § 4, no. 4). The same results are valid when ``topological group'' is replaced by ``topological ring''; we leave to the reader the task of stating the analogous results for topological modules (§ 6, no. 6).

Let \((\mathbf{X}_{\alpha}, g_{\alpha \beta})\) be an inverse system of topological spaces, and let \((\mathbf{G}_{\alpha}, f_{\alpha \beta})\) be an inverse system of topological groups. Suppose that each \(\mathbf{G}_{\alpha}\) operates continuously on \(\mathbf{X}_{\alpha}\) (§ 2, no. 4) and that the relations (1) hold for all \(x_{\beta} \in \mathbf{X}_{\beta}\) , \(s_{\beta} \in \mathbf{G}_{\beta}\) and \(\alpha \leqslant \beta\) . We have seen (no. 1) that \(\mathbf{X} = \varprojlim \mathbf{X}_{\alpha}\) has \(\mathbf{G} = \varprojlim \mathbf{G}_{\alpha}\) as group of operators; furthermore, \(\mathbf{G}\) operates continuously on \(\mathbf{X}\) . For if \(g_{\alpha}\) (resp. \(f_{\alpha}\) ) is the canonical mapping \(\mathbf{X} \to \mathbf{X}_{\alpha}\) (resp. \(\mathbf{G} \to \mathbf{G}_{\alpha}\) ), then by definition

\[
g _ {\alpha} (s, x) = f _ {\alpha} (s). g _ {\alpha} (x)
\]

and therefore the mappings \((s.x) \to g_{\alpha}(s.x)\) are continuous on \(\mathbf{X} \times \mathbf{G}\) , which shows the continuity of \((s, x) \to (s.x)\) (Chapter I, § 4, no. 4).

The mapping \(h_{\alpha}: \mathbf{X} / \mathbf{G} \to \mathbf{X}_{\alpha} / \mathbf{G}_{\alpha}\) induced by \(f_{\alpha}\) and \(g_{\alpha}\) is therefore continuous (§ 2, no. 4), and so is the mapping \(h: \mathbf{X} / \mathbf{G} \to \varprojlim \mathbf{X}_{\alpha} / \mathbf{G}_{\alpha}\) defined by the \(h_{\alpha}\) (Chapter I, § 2, no. 3, Proposition 4).

PROPOSITION 1. Suppose that the \(\mathbf{X}_{\alpha}\) and the \(\mathbf{G}_{\alpha}\) satisfy the above hypotheses. a) If the stabilizer of each point of \(\mathbf{X}_{\alpha}\) is a compact subgroup of \(\mathbf{G}_{\alpha}\) for each \(\alpha \in \mathbf{I}\) , then the stabilizer of each point \(x = (x_{\alpha})\) of \(\mathbf{X}\) is a compact subgroup of \(\mathbf{G}\) ; the orbit of \(x\) (with respect to \(\mathbf{G}\) ) is canonically homeomorphic to the inverse limit of the orbits of the \(x_{\alpha}\) (with respect to \(\mathbf{G}_{\alpha}\) ); and the canonical mapping \(h: \mathbf{X} / \mathbf{G} \to \varprojlim \mathbf{X}_{\alpha} / \mathbf{G}_{\alpha}\) is injective.

\begin{enumerate}
\def\labelenumi{\alph{enumi})}
\setcounter{enumi}{1}
\tightlist
\item
  If, for each \(\alpha \in \mathbf{I}\) , every orbit of a point of \(\mathbf{X}_{\alpha}\) (with respect to \(\mathbf{G}_{\alpha}\) ) is compact, then every orbit of a point of \(\mathbf{X}\) (with respect to \(\mathbf{G}\) ) is relatively compact, and \(h\) is surjective. If also \(h\) is bijective, then every orbit of a point of \(\mathbf{X}\) is compact.
\end{enumerate}

Let \(x = (x_{\alpha}) \in \mathbf{X}\) , and for each \(\alpha \in \mathbf{I}\) let \(\mathbf{X}_{\alpha}' = \mathbf{G}_{\alpha}.x_{\alpha}\) be the orbit of \(x_{\alpha}\) . If \(\alpha \leqslant \beta\) , it follows from (1) and the relation \(g_{\alpha\beta}(x_{\beta}) = x_{\alpha}\) that \(g_{\alpha\beta}(\mathbf{X}_{\beta}') \subset \mathbf{X}_{\alpha}',\) and hence that \((\mathbf{X}_{\alpha}')\) is an inverse system of subsets of the \(\mathbf{X}_{\alpha}\) . For each \(\alpha \in \mathbf{I}\) let \(u_{\alpha}: \mathbf{G}_{\alpha} \to \mathbf{X}_{\alpha}'\) be the continuous mapping \(s_{\alpha} \to s_{\alpha}.x_{\alpha}\) ; the \(u_{\alpha}\) form an inverse system of mappings, and \(u = \varprojlim u_{\alpha}\) is the continuous mapping \(s \to s.x\) of \(G\) into the subspace \(X' = \varprojlim X_{\alpha}'\) of \(X\) . The hypothesis of a) implies that \(\overline{u}_{\alpha}^{-1}(y_{\alpha})\) is compact for each \(y_{\alpha} \in X_{\alpha}'\) . Since also \(u_{\alpha}\) is surjective, the conditions of Chapter I, § 9, no. 6, Proposition 8, Corollary 2 are satisfied, and the first two assertions of a) are therefore established. Hence, if \(x = (x_{\alpha})\) and \(y = (y_{\alpha})\) are such that \(x_{\alpha}\) and \(y_{\alpha}\) belong to the same orbit with respect to \(G_{\alpha}\) for each \(\alpha \in I\) , then \(x\) and \(y\) belong to the same orbit with respect to \(G\) , and therefore \(h\) is injective.

Likewise, the hypothesis of b) implies that the inverse system of canonical mappings \(v_{\alpha}: \mathbf{X}_{\alpha} \to \mathbf{X}_{\alpha}/\mathbf{G}_{\alpha}\) satisfies the conditions of Chapter I, § 9, no. 6, Proposition 8, Corollary 2, and therefore its inverse limit \(v = \varprojlim v_{\alpha}: \mathbf{X} \to \varprojlim \mathbf{X}_{\alpha}/\mathbf{G}_{\alpha}\) is surjective, and the inverse image under \(v\) of every point of \(\varprojlim X_{\alpha}/G_{\alpha}\) is compact. Since \(v\) factorizes into

\[
\mathrm{X} \xrightarrow {\psi} \mathrm{X} / \mathrm{G} \xrightarrow {h} \varprojlim \mathrm{X} _ {\alpha} / \mathrm{G} _ {\alpha},
\]

where \(\psi\) is the canonical mapping, the assertions of b) follow.

COROLLARY 1. If the \(\mathbf{G}_{\alpha}\) are compact and the \(\mathbf{X}_{\alpha}\) Hausdorff, then the conclusions of a) and b) are valid.

For the hypotheses of a) and b) are satisfied, since every closed subgroup of \(G_{\alpha}\) is compact and \(u_{\alpha}: s_{\alpha} \to s_{\alpha}.x_{\alpha}\) is a continuous mapping of a compact space into a Hausdorff space.

COROLLARY 2. If, for each \(\alpha \in I\) , the group \(G_{\alpha}\) operates transitively on the space \(X_{\alpha}\) , and if the stabilizer of each point of \(X_{\alpha}\) is a compact subgroup of \(G_{\alpha}\) , then \(G\) operates transitively on \(X\) and the stabilizer of each point of \(X\) is a compact subgroup of \(G\) .

For hypothesis \(a)\) is satisfied, and \(\mathbf{X}_{\alpha}^{\prime} = \mathbf{X}_{\alpha}\) for each \(\alpha\) .

COROLLARY 3. Suppose that the \(\mathbf{G}_{\alpha}\) are Hausdorff. For each \(\alpha \in \mathbf{I}\) , let \(\mathbf{K}_{\alpha}\) be a compact subgroup of \(\mathbf{G}_{\alpha}\) , such that \(f_{\alpha \beta}(\mathbf{K}_{\beta}) \subset \mathbf{K}_{\alpha}\) whenever \(\alpha \leqslant \beta\) . Then, if \(\mathbf{K} = \varprojlim \mathbf{K}_{\alpha}\) , the canonical mapping \(h\) of the homogeneous space \(\mathbf{G} / \mathbf{K}\) into \(\varprojlim \mathbf{G}_{\alpha} / \mathbf{K}_{\alpha}\) is a homeomorphism.

The fact that \(h\) is bijective follows from Corollary 1, applied by replacing \(\mathbf{X}_{\alpha}\) by \(\mathbf{G}_{\alpha}\) , and \(\mathbf{G}_{\alpha}\) by \(\mathbf{K}_{\alpha}\) operating by right translations (§ 2, no. 5). Let \(\varphi\) be the canonical mapping \(\mathbf{G} \to \mathbf{G} / \mathbf{K}\) , and for each \(\alpha\) let \(f_{\alpha}\) be the canonical mapping \(\mathbf{G} \to \mathbf{G}_{\alpha}\) . If, for each \(\alpha, \mathbf{V}_{\alpha}\) runs through a fundamental system of open neighbourhoods of the identity element \(e_{\alpha}\) of \(\mathbf{G}_{\alpha}\) , then the sets \(\mathbf{V} = \overline{f}_{\alpha}^{-1}(\mathbf{V}_{\alpha}) (\alpha \text{ and } \mathbf{V}_{\alpha} \text{ variable})\) form a fundamental system of neighbourhoods of the identity element \(e\) in \(\mathbf{G}\) (Chapter I, § 4, no. 4, Proposition 9), and the sets \(\varphi(\mathbf{V}.K)\) form a fundamental system of neighbourhoods of \(\varphi(e)\) in \(G/K\) . We have to show that the image of \(\varphi(V.K)\) under \(h\) contains a neighbourhood of \(h(\varphi(e))\) , that is that there exists \(\beta \geqslant \alpha\) and a neighbourhood \(W_{\beta}\) of \(e_{\beta}\) in \(G_{\beta}\) such that \(\overline{f}_{\beta}^{-1}(W_{\beta}.K_{\beta}) \subset V.K\) . Now, the relation \(x \in V.K\) is equivalent to the existence of \(y\) in \(K\) such that \(f_{\alpha}(xy^{-1}) \in V_{\alpha}\) , i.e.~\(f_{\alpha}(x) \in V_{\alpha}.f_{\alpha}(K)\) ; so that \(V.K = \overline{f}_{\alpha}^{-1}(V_{\alpha}.f_{\alpha}(K))\) . Let \(U_{\alpha} = V_{\alpha}.f_{\alpha}(K)\) ; we shall see that there exists \(\beta \geqslant \alpha\) such that if we put \(U_{\beta} = \overline{f}_{\alpha\beta}^{-1}(U_{\alpha})\) , we have \(K_{\beta} \subset U_{\beta}\) ; it will then follow that there is a neighbourhood \(W_{\beta}\) of \(e_{\beta}\) in \(G_{\beta}\) such that \(W_{\beta}.K_{\beta} \subset U_{\beta}\) (Chapter II, § 4, no. 3, Corollary to Proposition 4), and this will establish the desired relation

\[
\overline {{f}} _ {\beta} ^ {-1} (\mathrm{W} _ {\beta}. \mathrm{K} _ {\beta}) \subset \overline {{f}} _ {\beta} ^ {-1} (\mathrm{U} _ {\beta}) = \mathrm{V}. \mathrm{K}.
\]

We proceed by reductio ad absurdum: for each \(\beta \geqslant \alpha\) , let \(\mathbf{M}_{\beta}\) denote \(\mathbf{K}_{\beta} \cap \mathbb{C}\mathbf{U}_{\beta}\) ; since \(\overline{f}_{\beta \gamma}^{-1}(\mathrm{U}_{\beta}) = \mathrm{U}_{\gamma}\) if \(\alpha \leqslant \beta \leqslant \gamma\) , the \(\mathbf{M}_{\beta}\) form an inverse system of compact subsets of the \(\mathbf{G}_{\beta}\) (for \(\beta \geqslant \alpha\) ). If they were all non-empty, then their inverse limit \(\mathbf{M}\) would also be non-empty (Chapter I, § 9, no. 6, Proposition 8). It is clear that \(\mathbf{M} \subset \mathbf{K}\) and \(f_{\alpha}(\mathbf{M}) \subset \mathbf{M}_{\alpha}\) ; but this is absurd since \(f_{\alpha}(\mathbf{K}) \subset \mathbf{U}_{\alpha}\) , and the proof is therefore complete.

